\section{An independent classical bound for quadratic cores}\label{sec:older-route}

This section proves the order-of-magnitude input for the quadratic-order
application without the zero-free half-plane or weighted-dilation
theorem.  It is the quadratic specialization of the divisor-pairing
and positive-restoration argument.  Here the first divisor product has
only two factors.  Consequently ordinary completed Kloosterman sums
supply the required progression estimate; no higher divisor
correlation theorem is needed.

Fix an even modulus \(L_0\), a nonprincipal quadratic literal character
\(\chi\bmod L_0\), and put
\[
 \begin{gathered}
 H=\ker\chi,\qquad
 \mathcal C=\langle p:p\nmid L_0,\ \chi(p)=1\rangle,\\
 V(n)=\tau(n)\1_{\mathcal C}(n),\qquad
 J(N)=\prod_{\substack{p\mid N\\p\nmid L_0,\ \chi(p)=-1}}
                          \frac{p-1}{p-2}.
 \end{gathered}
\]
All arithmetic data in the following lemma are fixed.

\begin{lemma}[Classical quadratic core bound]\label{lem:older-quadratic-filter}
Let \(A,B\) be coprime positive \(L_0\)-supported integers and
\(U,V_0\subset H\) fixed residue subsets.  Write
\[
 \Gamma_{A,B;U,V_0}(N)=\frac{L_0}{|H|^2}
   \#\{(a,b)\in U\times V_0:Aa+Bb=N\bmod L_0\}.
\]
Uniformly over all sufficiently large targets with \(\Gamma_{A,B;U,V_0}(N)>0\),
\[
 \#\{Ax+By=N:x,y\in\mathcal C,\
                  x\bmod L_0\in U,\ y\bmod L_0\in V_0\}
                   \asymp\frac{NJ(N)}{\log N}.
\]
Zero finite density gives zero count identically.  In particular the
lemma applies to unions of unit-square residue states.
Its proof uses only fixed-character means, fixed-progression prime
Mertens, the classical Weil--Estermann bound, and finite sieve inequalities.
\end{lemma}

We prove the full shifted weighted family first.  Its uniformity,
rather than a single weighted convolution, is necessary for restoring
small prime factors.

\subsection{Two-factor boxes with a fixed period}

For bounded functions \(\omega_1,\omega_2\) of fixed period \(c\),
intervals \(I_1,I_2\), and \((a,m)=(c,m)=1\), let
\[
 S(a)=\sum_{\substack{r\in I_1,\ s\in I_2\\rs\equiv a\ (m)}}
                        \omega_1(r)\omega_2(s).
\]
Then, for every \(\epsilon>0\),
\begin{equation}\label{eq:older-box}
 S(a)-\frac1{\varphi(m)}\sum_{b\bmod m}^{*}S(b)
                         \ll_{c,\epsilon}m^{1/2+\epsilon}.
\end{equation}
The estimate is uniform in the intervals, including intervals longer
than \(cm\), and permits subintervals of arbitrary lengths.

Here is a direct completion proof.  Put \(Q=cm\),
\[
 W_i(h)=\sum_{n\in I_i}\ee{-hn/Q},\qquad
 G_i(t)=\sum_{r\bmod c}\omega_i(r)\ee{tr/c}.
\]
CRT and finite Fourier inversion express \(S(a)\) as
\[
 Q^{-2}\sum_{h,k\bmod Q}
 W_1(h)W_2(k)G_1(m^{-1}h)G_2(m^{-1}k)
       K_1(c^{-1}h,ac^{-1}k;m),
\]
where \(K_1(u,v;m)=\sum_{z\bmod m}^{*}\ee{(uz+v\bar z)/m}\);
inverses in \(G_i\) are taken modulo \(c\).
The terms \(h=0\bmod m\) or \(k=0\bmod m\) are independent of
the unit \(a\), since the resulting Ramanujan sum is unchanged by
unit multiplication.  For all other frequencies use
\[
 |K_1(u,v;m)|\le\tau(m)m^{1/2}(u,v,m)^{1/2}.
\]
This is the ordinary Weil--Estermann bound stated and applied in
\cite[printed p.~44]{heathbrown}.  The primary source's
Lemma~6 is its singleton-padded divisor-box consequence.
In centered frequency ranges,
\(|W_i(h)|\ll\min(|I_i|,Q/|h|)\) and \(|G_i|\le c\).
The remaining double sum is bounded using
\[
 \sum_{\substack{0<|h|,|k|\le Q/2\\m\nmid h,\ m\nmid k}}
      \frac{(h,k,m)^{1/2}}{|hk|}
 \le \sum_{d\mid m}d^{1/2}
       \left(\sum_{\substack{0<|h|\le Q/2\\d\mid h}}\frac1{|h|}\right)^2
 \ll_c \log^2(2m)\sum_{d\mid m}d^{-3/2}.
\]
The \(Q^2\) from the two geometric sums cancels the completion
normalization.  This proves a bound
\(O_c(m^{1/2}\tau(m)\log^2(2m))\) around an \(a\)-independent
quantity.  Averaging that quantity over unit \(a\) changes the
error by at most a factor two and proves \eqref{eq:older-box}.
The case \(m=1\) has discrepancy zero.
Tensor partial summation also allows bounded-variation interval weights.

\subsection{Shifted kernels, means, and fixed filters}

For \(t\in\mathcal C\) put \(F_t(n)=V(tn)/V(t)\).
At a split prime with \(b=v_p(t)\),
\[
 F_t(p^e)=1+\frac{e}{b+1};
\]
at a nonsplit prime it is zero for \(e>0\).
With literal characters, its Dirichlet series is
\[
 \sum_{n\ge1}\frac{F_t(n)}{n^s}
   =L(s,\chi_0^{L_0})L(s,\chi)\sum_{j\ge1}\frac{e_t(j)}{j^s}.
\]
The correction factors are \(1-p^{-2s}\) at inert primes and
\(1-\frac{b}{b+1}p^{-s}\) at split primes dividing \(t\).
Deleting all Euler factors at primes dividing another integer \(D\)
replaces its split local correction by \((1-p^{-s})^2\).
Thus all these kernels satisfy, uniformly in valuations,
\begin{equation}\label{eq:older-kernel}
 |e(j)|\le\tau(j)^C,\qquad
 \sum_j|e(j)|j^{-3/5}\ll_{\epsilon}(tD)^\epsilon.
\end{equation}
Every kernel index is an \(L_0\)-unit with \(\chi(j)=1\).

We record explicitly the elementary mean used below.  If
\(\xi\) is a character of a fixed modulus \(c\) with the same prime
support as \(L_0\), then the twisted series has character factors
\(\xi,\xi\chi\).
It has a principal factor only when the restriction of \(\xi\)
is \(1\) or \(\chi\); in those cases \(F_t\xi=F_t\).
Otherwise both character sums have bounded partial sums, and
Dirichlet hyperbola at \(\sqrt X\) gives \(O_c(\sqrt X)\) for
their convolution.
In the principal cases, split the principal factor from the
nonprincipal one by the same hyperbola argument to obtain its residue
times \(X\), with error \(O_c(\sqrt X)\).
Summing kernels by \eqref{eq:older-kernel}, including their main tails,
therefore gives
\begin{equation}\label{eq:older-mean}
 \sum_{n\le X}F_t(n)\1_{(n,D)=1}\xi(n)
 =\delta_{\xi\in\{1,\chi\}}\rho_{t,D}X
                         +O_\epsilon((tDX)^\epsilon X^{3/5}).
\end{equation}
Here the two displayed characters mean their inflations to \(c\).
The error and main remain valid with any fixed periodic unit filter,
after finite character expansion.
For later normalization write
\[
 M_{t,p}=(1-p^{-1})\sum_{e\ge0}F_t(p^e)p^{-e},\qquad
 \rho_t=\frac{\varphi(L_0)}{L_0}\prod_{p\nmid L_0}M_{t,p}.
\]
The product is taken in natural prime order; its convergence and
positivity follow equally from the character factorization and
\(L(1,\chi)>0\).  A deleted split prime divides this mean by its
original local series.  The residue measure on the fixed unit
group is uniform on \(H\).

\begin{proposition}[Uniform shifted weighted family]\label{prop:older-weighted}
There exist \(\eta,\theta>0\), with \(\theta\) chosen sufficiently
small after \(\eta\), such that the following holds.
Let \(t_1,t_2,u,v\in\mathcal C\), each at most \(N^\theta\), with
no restriction on their gcds.  For fixed residue subsets \(U,V_0\subset H\),
\begin{align*}
 &\sum_{\substack{Au x+Bv y=N\\x\bmod L_0\in U,\ y\bmod L_0\in V_0}}
                      F_{t_1}(x)F_{t_2}(y)\\
 &\quad=\frac{N}{ABuv}\,
       \Gamma_{Au,Bv;U,V_0}(N)\rho_{t_1}\rho_{t_2}
       \prod_{p\nmid L_0}\frac{D_p}{M_{t_1,p}M_{t_2,p}}
                         +O(N^{1-\eta}),
\end{align*}
uniformly in every target valuation.  Here
\[
 D_p=\lim_{h\to\infty}p^{-h}
 \sum_{\substack{a,b\bmod p^h\\Au a+Bv b=N\ (p^h)}}
       F_{t_1}(p^{\min(v_p(a),h)})
       F_{t_2}(p^{\min(v_p(b),h)}).
\]
All factors, including zero factors, are retained.
The same formula holds after imposing any fixed finite combination
of small divisibility and exact-valuation tests whose resulting
shifts and coefficients stay in this family.
\end{proposition}

\begin{proof}
First suppose \((u,v)=(uv,N)=1\).
For \(g=(x,N)\), write \(x=gx'\), \(y=gy'\), \(M=N/g\);
a nonzero term has \(g\in\mathcal C\).
Its exact prefactor is \(F_{t_1}(g)F_{t_2}(g)\), its new shifts
are \(t_1g,t_2g\), and \((x',M)=1\).
Transport the residue states to \(U/g,V_0/g\).
Every such multiplication is a unit bijection, so the corresponding
finite root count, with its transported target, is unchanged.

Choose a small fixed \(\lambda>0\), put \(K=N^\lambda\), and truncate
\(g\) and the two kernels at \(K\).
Large gcds cost \(O(N^{1-\lambda+\epsilon})\).
For either kernel, counting multiples and using
\eqref{eq:older-kernel} gives
\[
 N^{1+\epsilon}\sum_{j>K}\frac{|e(j)|}{j}
   +N^\epsilon\sum_{j\le N}|e(j)|
       \ll N^{1-2\lambda/5+\epsilon}+N^{3/5+\epsilon}.
\]
Discard also \(Au x\le N^{1-\lambda}\), at
\(O(N^{1-\lambda+\epsilon})\).
All implicit valuation and shift factors here are divisor-bounded.

We suppress primes on \(x',y'\) in the remainder of this core argument.
The fixed filter on the original first unit variable is
\[
 w(x)=\1_{x\bmod L_0\in U/g}\,
       \1_{B\mid M-Au x}\,
       \1_{(z,L_0)=1,\ z\bmod L_0\in v(V_0/g)},\qquad
                       z=(M-Au x)/B.
\]
It has fixed period \(c=BL_0\).  Expand it on unit classes by characters
modulo \(c\), retaining every scalar phase of every kernel index.
The filter ensures that
\(t=(M-Au x)/(Bvk)\) is an \(L_0\)-unit with \(\chi(t)=1\).
Consequently the exact symmetric divisor pairing is
\[
 (\chi_0^{L_0}*\chi)(t)
  =2\sum_{\substack{a\mid t\\a<\sqrt t}}\chi(a)
                         +\1_{t=\square}\chi(\sqrt t).
\]
The square endpoint costs \(O(N^{1/2+O(\lambda+\theta)+\epsilon})\).

For a second kernel \(k\le K\), the remaining congruence is
\[
       Au x\equiv M\pmod{vka},\qquad
       Au x+Bvka^2<M,\qquad (x,M)=1 .
\]
The entire bad-supported factor \(B\) has already been imposed in
\(w\); all varying moduli are coprime to \(c\).
Also \((k,uM)=(a,uM)=1\).
Expand the first kernel \(j\) and its two character variables.
A nonzero congruence requires \((uj,vka)=1\); otherwise it would
force a prime divisor of \(uM\) into an excluded place, or make a
unit equal to a nonunit.  The same zero test holds in the unit model.

Partition the two character variables and \(a\) into relative-width
\(\Delta=N^{-\zeta}\) boxes, with \(\zeta=\lambda/100\).
There are \(O(N^{3\zeta+\epsilon})\) boxes.
On a box with its endpoints fixed, apply \eqref{eq:older-box} to
the first two factors, at the full modulus \(m=vka\).
It is enough that \(m\le N^{1/2+O(\lambda+\theta)}\);
summing its square-root error over \(a\le N^{1/2}\),
then the two kernel masses, gives
\begin{equation}\label{eq:older-box-budget}
 O(N^{3/4+100\lambda+100\theta+\epsilon}).
\end{equation}
This deliberately loose bound includes the fine boxes, the fixed
character expansions, and the \(g\)-strata.
In particular, no medium-variable dispersion argument occurs
in this quadratic specialization.

We justify the curved boundary separately.  A crossing box satisfies
\(|Au x+Bvka^2-M|\ll\Delta N\).
On the actual congruence write \(M-Au x=Bvkat\).
Then \(|Bvka(t-a)|\ll\Delta N\), which has
\[
 O\left(\frac{\Delta N\log N}{vk}+\sqrt{\frac N{vk}}\right)
\]
integer pairs \((a,t)\), after discarding only the divisibility
which determines \(x\).
The unique scalar decompositions and divisor-bounded first-kernel
multiplicities cost \(N^\epsilon\).
Summing the second kernel by \eqref{eq:older-kernel} leaves
\begin{equation}\label{eq:older-boundary}
 O_\epsilon\big(N^\epsilon(\Delta N\log N+\sqrt N K^{1/10})\big).
\end{equation}
For the principal model the interval length retains \(1/j\),
and kernel reciprocal sums converge.
After \eqref{eq:older-box}, evaluate its unit-population mean using
\eqref{eq:older-mean}, with Euler deletions at \(M vka\).
The \(1/\varphi(vka)\) normalization makes the sum of these mean
errors \(O(N^{3/5+O(\lambda+\theta)+\epsilon})\).
All fixed twists except \(1,\chi\) have zero mean.
Thus the remaining model can be evaluated directly, with a power saving.

For clarity, its finite coefficient can be checked before any Euler
completion.  Let \(H_c\) be the lift of \(H\) to modulus \(c=BL_0\).
It has \(B|H|\) elements.  Write \(\mu\) for the average of \(w\)
over \(H_c\).  There are exactly \(B\) lifts modulo \(BL_0\) of
each allowed pair modulo \(L_0\), since \(Au\) is a unit at primes
of \(B\); each surviving first residue admits \(B\) second lifts.
Therefore
\[
 \Gamma_{Au,Bv;U/g,V_0/g}(M)
                  =\frac{2BL_0}{\varphi(L_0)}\mu.
\]
The surviving first character coefficients sum to \(\mu\), rather
than to the constant coefficient alone.

Here are the radial factors in the model calculation.
For local coefficients \(A_i=F_{t_1g}(p^i)\),
\(B_i=F_{t_2g}(p^i)\), put \(q_p=1-p^{-1}\).
At a unit target, the local integral is
\(I_{p,0}=M_{A,p}+M_{B,p}-1\).
Writing \(W_p=M_{A,p}^{-1}\), the paired-divisor series contributes
\(I_{p,0}/M_{A,p}\) outside \(uvM\), \(1\) at \(p\mid u\),
and \(M_{B,p}/M_{A,p}\) at \(p\mid v\).
At \(p\mid M\), the first mean deletes its local series.
Hence the fixed-\(g\) radial product consists of
\[
 p\mid M:q_p,\qquad p\mid u:M_{A,p},\qquad
 p\mid v:M_{B,p},\qquad p\nmid uvM:I_{p,0},
\]
with real scale \(M/(Auv)\) and finite coefficient
\(2(\varphi(L_0)/L_0)\mu\).
These identities follow also by expanding the two geometric
series: the paired-divisor local sum off \(uvM\) is
\[
 1+W_p\left[
  \frac{\sum_i e_{t_2g}(p^i)p^{-i}}{1-\chi(p)/p}-1\right]
                  =\frac{M_{A,p}+M_{B,p}-1}{M_{A,p}} .
\]
Completing that character series costs
\(O(N^{1/2+O(\lambda+\theta)+\epsilon})\):
it is \(L(s,\chi)\) times a correction absolutely convergent
for \(\Re s>0\), so its coefficient partial sums are \(O(N^\epsilon)\).
The kernel completions have the saved tail already displayed.

Completing \(g\) gives, at \(n=v_p(N)\), exactly
\begin{align*}
 I_{p,n}={}&q_p\sum_{i<n}A_iB_ip^{-i}
        +p^{-n}(1-2/p)A_nB_n\\
 &+q_p\left(B_n\sum_{i>n}A_ip^{-i}
                       +A_n\sum_{j>n}B_jp^{-j}\right).
\end{align*}
The terms \(i<n\) come from the smaller common valuations;
the final common valuation \(i=n\) contributes its shifted
unit-target integral.  Thus no target valuation is omitted.
Since \(\rho_{t_1}\rho_{t_2}
=(\varphi(L_0)/L_0)^2\prod M_{t_1,p}M_{t_2,p}\), the
finite coefficient and real scale give precisely the formula
in the proposition, including its factor \(1/B\).

For general good coefficients, remove \(r=(u,v)\), requiring
\(r\mid N\), and divide them and the target by \(r\).
For coprime remaining coefficients put \(d=(u,N)\), \(e=(v,N)\).
The equation forces \(e\mid x,d\mid y\).
Substitute \(x=ex'\), \(y=dy'\), with exact prefactors
\(F_{t_1}(e)F_{t_2}(d)\), shifts \(t_1e,t_2d\), coefficients
\(u/d,v/e\), and target \(N/(de)\).
All removed factors are in \(\mathcal C\).
Transport each residue subset by the corresponding unit multiplier.
The finite solution sets are bijected by these substitutions.
At a common coefficient prime the local density gains \(p^{v_p(r)}\);
at a one-sided coefficient prime, solving for the unit-coefficient
variable gives the displayed \(D_p\).
This proves the claimed integral for every gcd profile.

Choose \(\lambda\) small, then \(\theta\) much smaller.
Equations \eqref{eq:older-box-budget} and
\eqref{eq:older-boundary}, the kernel saving \(2\lambda/5\),
and the gcd and small-summand saving \(\lambda\), leave a fixed
positive core saving.  Choose \(\eta\) below a hundredth of the
boundary saving, and decrease the external family exponent again.
The transformed targets are at least \(N^{1-3\theta}\) and all
transformed parameters stay inside the core family.
Divisor prefactors preserve a smaller saving.
Finally \(b\mid x\) is the exact substitution \(x=bx'\), which
changes the shift to \(t_1b\), coefficient to \(ub\), state to
\(U/b\), and introduces \(F_{t_1}(b)\).
The analogous substitution on the second side, and finite differences
for exact valuations, prove the predicate assertion.
\end{proof}

\subsection{A finite lower sieve and positive restoration}

We now prove Lemma~\ref{lem:older-quadratic-filter}.
Fix an admissible target.  Its finite density \(\Gamma(N)\)
belongs to a fixed finite positive set, and hence has a positive
lower bound depending only on the fixed data.
At good primes write \(x=p^{-1}\), \(q_p=1-x\), and use the
unshifted coefficients \(A_i=B_i=V(p^i)\).
Let \(I_{p,e}\) be the radial integral in the preceding proof,
where \(e=v_p(N)\).
The weighted mass with both variables units is
\[
 c_e(p)=
 \begin{cases}1-2/p,&e=0,\\1-1/p,&e\ge1.\end{cases}
\]
At an inert prime \(I_{p,e}=c_e(p)\).
At a split prime define the bad-event probability
\(g_e(p)=1-c_e(p)/I_{p,e}\).
The explicit geometric series gives, uniformly in \(e\),
\[
             0\le g_e(p)<1,\qquad
             g_e(p)=4/p+O(p^{-2}).
 \tag{A.1}\label{eq:older-roughprob}
\]
For example, at \(e=0\),
\(I_{p,0}=2(1-p^{-1})^{-1}-1\).
For \(e\ge1\), the \(i=0,1\) terms give
\(I_{p,e}=1+3/p+O(p^{-2})\); the remaining polynomial-geometric
tail is uniformly \(O(p^{-2})\).
The formula at \(e=1\) has the same expansion directly.
At every finite collection of primes \(I_{p,e}\) is uniformly bounded,
whereas \(c_e(p)\ge1/3\), so the good-event probabilities have a
uniform positive lower bound.

Put \(M_p=(1-p^{-1})\sum_iV(p^i)p^{-i}\) and
\[
 \mathcal V(z)=\prod_{\substack{p\le z\\p\nmid L_0}}
                  \frac{1-2/p}{M_p^2},\qquad
 E_p=\frac{p-1}{p-2},\qquad
 E_{p,e}=\begin{cases}1,&e=0,\\E_p,&e\ge1,\end{cases}
 \qquad a_p=\frac{p-1}{p(p-2)}.
\]
At a split prime \(M_p=(1-p^{-1})^{-1}\), so the factor
in \(\mathcal V\) is \(1-4/p+O(p^{-2})\).
At an inert prime it is \(1+O(p^{-2})\).
Fixed-character prime Mertens, and the split-prime density \(1/2\),
therefore give
\begin{equation}\label{eq:older-roughdimension}
                         \mathcal V(z)\asymp(\log z)^{-2}.
\end{equation}
This is the dimension-two rough weighted count which must recover
one power of \(\log N\).

Choose small fixed \(\alpha,\xi>0\) after the weighted family.
We use squarefree pairwise coprime triples \(d,u,v\) with
\[
 d,u,v\in\mathcal C,\quad d\mid N,\quad (uv,N)=1,\quad
 duv\le N^\alpha,
\]
all primes at most \(w=N^\omega\).
At target \(M=N/d\), prescribe \(v_p(x)=1\) for \(p\mid u\),
\(v_p(y)=1\) for \(p\mid v\), and require that no other prime
at most \(z=N^\gamma\) divide \(xy\).
Count with weight \(V(x)V(y)/(2^{\omega(u)}2^{\omega(v)})\)
and with residue states \(x\in U/d,\ y\in V_0/d\).
Thus the lifted values \(dx,dy\) have the original states.

All these predicates are covered by
Proposition~\ref{prop:older-weighted}.
Indeed exact valuation one is the difference of divisibility by
\(p\) and by \(p^2\), and a roughness predicate expands using
\[
 \1_{p\mid xy}=\1_{p\mid x}+\1_{p\mid y}
                              -\1_{p\mid x,\ p\mid y}.
\]
For a tested squarefree divisor \(\ell\le N^\xi\), each resulting
shift and coefficient is at most \(N^{2\alpha+\xi}\).
The expansion has \(N^{o(1)}\) terms.
It suffices to choose
\begin{equation}\label{eq:older-smallbudgets}
 2\alpha+\xi<(1-\alpha)\theta,\qquad
                         \alpha+\xi<\eta/4 .
\end{equation}
Forbidden shifts give zero terms; common-divisibility terms use the
noncoprime part of the weighted family.

The finite density remains exactly \(\Gamma(N)\) in every lower term.
After writing \(x=ux'\), \(y=vy'\), its states are
\(U/(du),V_0/(dv)\) and its equation is
\(Au x'+Bv y'=N/d\).
The map
\((x',y')\mapsto(du x',dv y')\) bijects its finite roots with the
original roots modulo \(L_0\).
This is why split primes with nontrivial residue states may be
restored.  Restricting restoration to identity-state primes would
lose the exponent needed below.

Here is the exact positive local model after imposing all roughness
conditions.  Its main mass for a triple is
\begin{equation}\label{eq:older-triplemass}
 \mathcal M_{d,u,v}
 =C_N(z)\frac{N}{ABd}\mathcal V(z)
       \prod_{\substack{p\mid N/d\\p\le z,\ p\nmid L_0}}E_p
       \prod_{p\mid uv}a_p .
\end{equation}
The positive factor \(C_N(z)\) includes \(\Gamma(N)\rho_1^2\)
and the Euler tail above \(z\).
It is bounded above and below by fixed positive constants,
uniformly in the target and the triples.
To verify this, outside the target the tail factors are
\(1+O(p^{-2})\); at target primes they are \(1+O(1/p)\) uniformly
in valuation.  At most \(1/\gamma\) target primes exceed \(z\).
All cofactor primes are below \(w<z\), so this tail is unchanged
when dividing by \(d\).
At a unit target prime, a valuation-one directive on either side
has weighted mass \(2(1-p^{-1})/p\).
Its normalization by two and division by \(1-2/p\) give exactly
\(a_p\), proving \eqref{eq:older-triplemass}.

We justify that a finite sieve recovers a fixed fraction of this
fully rough main, rather than asserting an asymptotic sieve.
Treat finitely many small primes exactly so that \(g_e(p)\le1/2\)
at the others.  Split the remaining primes into blocks
\((z^{2^{-j-1}},z^{2^{-j}}]\).
Their sums of local bad probabilities are bounded by a fixed
constant \(M_0\), uniformly in targets and directives.
Let \(U_j\) be the even Bonferroni polynomial through degree
\(2(s+j)\), and \(L_j\) the odd one through degree \(2(s+j)+1\).
If \(I_j\) is the indicator of no bad event in block \(j\), then
\(U_j\ge I_j\ge L_j\), \(U_j\ge0\), and telescoping gives
\begin{equation}\label{eq:older-lowerpoly}
 \prod_jU_j-\sum_j(U_j-L_j)\prod_{i\ne j}U_i
                                      \le\prod_jI_j .
\end{equation}
This avoids multiplying lower polynomials which might be negative.
In the independent prime model,
\[
 \E(U_j-L_j)\le\frac{M_0^{2(s+j)+1}}{(2(s+j)+1)!},
 \qquad \E U_j\ge\prod_{p\text{ in block}}(1-g_e(p))
                                             \ge e^{-2M_0}.
\]
Choose the fixed \(s\) so that the sum of these relative errors
is at most \(1/4\).
Factoring \(\prod_j\E U_j\) shows that
\eqref{eq:older-lowerpoly} has mean at least half the fully rough
probability.  Its squarefree divisor support is at most
\[
 C_0z^{\sum_{j\ge0}(2(s+j)+1)2^{-j}}
                               =C_0z^{4s+6},
\]
with coefficients \(O(\log\log N)\).
Choose \(\gamma<\xi/(4s+7)\).
Then its support is at most \(N^\xi\), so
Proposition~\ref{prop:older-weighted} applies term by term.
There are at most \(N^{\alpha+o(1)}\) cofactor triples.
The total error, including sieve predicates, is
\[
                  O(N^{1-\eta+\alpha+\xi+o(1)}),
\]
negligible by \eqref{eq:older-smallbudgets}.
Thus the summed actual weighted count is at least a fixed fraction
of the sum of \eqref{eq:older-triplemass}, apart from that error.

Temporarily drop the size cap \(duv\le N^\alpha\).
The positive cofactor sum is an Euler product.
At a split prime not dividing \(N\), its options are
\(1,a_p,a_p\).
At a split prime with target valuation \(e\ge1\), the options
omitting or including it in \(d\) give
\[
 E_p+p^{-1}E_{p,e-1}\ge E_p+1/p=1+2a_p.
\]
At inert target primes no restoration option is present; the
factor left over is \(E_p\).
Target factors above \(z\) contribute \(1+o(1)\), uniformly.
Consequently \eqref{eq:older-roughdimension} and prime Mertens give
\begin{equation}\label{eq:older-restoredmass}
 \sum_{d,u,v}\mathcal M_{d,u,v}
 \gg NJ(N)(\log N)^{-2}
          \prod_{\substack{p\le w\\\chi(p)=1}}(1+2a_p)
 \asymp \frac{NJ(N)}{\log N}.
\end{equation}

The cap must be restored in the actual positive model.
After normalizing its Euler product, each prime is included in
\(d,u\), or \(v\) with probability \(O(1/p)\), uniformly in target
valuation.  Therefore
\(\E\log(duv)\le C_1\log w\).
Choose \(\omega<\min(\gamma,\alpha/(4C_1))\).
Markov retains at least three quarters of the mass in
\(duv\le N^\alpha\).  Applied to this model, rather than merely a
comparison product, it proves that the finite sieve retains the
order of magnitude in \eqref{eq:older-restoredmass}.

Finally lift each residual pair to \((dx,dy)\).
At every cofactor prime the numerical valuation patterns are,
respectively,
\((1,1)\) for \(d\), \((1,0)\) for \(u\), and \((0,1)\) for \(v\);
all other small-prime valuations are zero.
Thus the ordered numerical pair recovers the entire triple and
residual pair uniquely.
Writing \(x=ux'\), \(y=vy'\), exact valuation one gives
\[
       \frac{V(x)V(y)}{2^{\omega(u)}2^{\omega(v)}}
                                      =V(x')V(y').
\]
All prime factors of \(x',y'\) exceed \(N^\gamma\), so
\(V(x')V(y')\le4^{1/\gamma}\).
No factor \(V(d)^2\) is inserted: \(d\) lifts values, not their
divisor weights.
Injectivity and this uniform bound turn
\eqref{eq:older-restoredmass} into the required unweighted lower bound.

For the upper bound, sieve the integer parameter of \(Ax+By=N\)
using only good inert primes.  Both coefficients are units there.
There are two forbidden roots, merging exactly when \(p\mid N\).
A dimension-one upper Selberg sieve at fixed power level therefore
gives
\[
 \ll N\prod_{\substack{p\le N^{\gamma_1}\\
                       p\nmid L_0,\ \chi(p)=-1}}
       \left(1-\frac{2-\1_{p\mid N}}p\right)
            +O(N^{1/2}\log^C N)
 \ll \frac{NJ(N)}{\log N}.
\]
For an explicit finite version, use the nonnegative even block
polynomials above at divisor support at most \(N^{1/2}\).
Their model is at most a fixed multiple of the true product,
by the same summable factorial block errors.
CRT gives \(N\nu(\ell)/\ell+O(\nu(\ell))\), \(\nu(p)\le2\);
the total endpoint error is
\(O(N^{1/2}\log^C N)\).
Discarding the residue filters only increases this upper bound.
This completes the proof of Lemma~\ref{lem:older-quadratic-filter}.

\paragraph{Independence of the two conclusions.}
Lemma~\ref{lem:older-quadratic-filter} supplies uniform two-sided
bounds, but no leading asymptotic.  Its positivity comes from the
finite lower sieve and the recoverable positive cofactor model.
The density-model comparison of the main binary theorem supplies
the leading asymptotic by a separate argument.  In particular the
classical bound proved here has no dependency on the recent
zero-free or graph inputs.
